\section{背景}

\subsection{消费级设备特征}

我们关注的消费级设备通常具有以下特征：
\begin{itemize}
\item \textbf{资源受限}：主流笔记本通常只有 8GB 显存和 16--32GB 内存，远低于数据中心服务器。
\item \textbf{单机部署}：无需分布式协调，但要求系统足够轻量且高效。
\item \textbf{稳定性要求高}：在家庭或工业场景中，系统崩溃会显著损害用户体验。
\end{itemize}

\subsection{多模态任务特征}

我们使用三个正交属性来刻画消费级本地 AI 运行时中的工作负载：服务语义、设备亲和性和时延敏感性。相较于只按模型家族分类，这种视角更能直接反映本地运行时中的冲突。

\textbf{服务语义}：
\begin{itemize}
\item \textbf{常驻任务}：长期运行、应持续可用、不能容忍长时间无响应的服务。交互式 LLM 助手是典型例子。
\item \textbf{一次性任务}：以独立请求形式提交、执行并完成的突发任务。图像生成是本文中的代表性实例。
\end{itemize}

\textbf{设备亲和性}：
\begin{itemize}
\item \textbf{GPU 绑定任务}：若不在独占 GPU 上执行，其时延将不可接受。在消费级笔记本上，扩散式图像生成通常属于此类。
\item \textbf{可迁移任务}：虽然在 GPU 上性能更好，但在合适的量化和推理栈支持下，退化到 CPU 仍可继续提供可用服务。量化 LLM 服务就是本文中的代表。
\end{itemize}

\textbf{时延敏感性}：
\begin{itemize}
\item \textbf{交互式任务}：直接面向用户，需要较低延迟和持续响应。
\item \textbf{可排队任务}：可以容忍适度等待，只要最终能够正确完成即可。
\end{itemize}

在这一抽象下，当前的 EdgeSched 实例聚焦于本地 AI 运行时中一对冲突尤为突出的工作负载：LLM 服务作为常驻、可迁移、时延敏感的服务，图像生成作为突发、GPU 绑定、可排队的任务。下文都将以这一对工作负载作为调度策略呈现与评估的具体实例。

\subsection{设计目标}

基于上述工作负载抽象，EdgeSched 的设计目标包括：

\begin{enumerate}
\item \textbf{轻量化}：调度器本身必须保持轻量（启动快、内存占用小），而模型自身的资源使用应单独核算。
\item \textbf{降级条件下的服务连续性}：在 GPU 竞争期间，常驻的交互式任务即便退化到 CPU 执行，也应保持可用。
\item \textbf{高稳定性}：系统应能够在崩溃后快速恢复到语义上可用的状态，避免任务状态丢失或悬挂。
\item \textbf{运维务实性}：系统应便于部署、调试和复现实验结果，适合单机本地环境。
\end{enumerate}